
\chapter{GIỚI THIỆU TỔNG QUAN}

\section{Bối cảnh: Sự phát triển của Tương tác Cộng tác Người - Máy (HRC)}
Trong những thập kỷ gần đây, ngành công nghiệp tự động hóa đã chứng kiến những bước tiến đáng kể trong lĩnh vực vận hành robot. Trước kia, nhu cầu sử dụng robot chỉ xuất hiện để xử lý những công việc có tính chất lặp đi lặp lại, hoặc cần thực hiện chính xác theo một chuỗi chương trình đã được lập trình sẵn. Ngày nay, chức năng của robot không chỉ giới hạn ở việc hoạt động độc lập bên trong các lồng kính an toàn tại các dây chuyền sản xuất hàng loạt, mà đang dần chuyển dịch sang mô hình Tương tác cộng tác Người - Máy (Human-Robot Collaboration - HRC) \cite{wang2024genealogy}. Mục tiêu của sự dịch chuyển này là kết hợp sự bền bỉ, tính chính xác của máy móc với khả năng nhận thức, ra quyết định linh hoạt của con người trước những tình huống phức tạp, phi cấu trúc mà không thể lập trình trước.

Sự cộng tác này đặc biệt cần thiết trong các môi trường làm việc đặc thù như sản xuất linh kiện phức tạp, y tế, xử lý vật liệu độc hại, hay không gian vũ trụ \cite{pryor2023virtualspace}. Tính chất của công việc ở các môi trường này đều có điểm chung: tiềm ẩn rủi ro về an toàn cho con người. Chính vì vậy, tương tác Người - Máy được áp dụng trong các môi trường này sẽ chủ yếu là Teleoperation (điều khiển từ xa), khi đó người vận hành sẽ sử dụng các thiết bị đặc trưng để vận hành thiết bị ở môi trường rủi ro, qua đó đảm bảo an toàn cho người vận hành, mặt khác giúp cho thiết bị có thể thực hiện những công việc phức tạp cần có kĩ năng của con người. Một ví dụ trong lĩnh vực y tế, đặc biệt trong thời kì covid, việc kiểm tra sức khỏe bệnh nhân tại những nơi có tỉ lệ nhiễm bệnh cao tạo rủi ro nhiễm bệnh cho chính các bác sĩ, khiến cho lợi ích của việc ứng dụng robot trong lĩnh vực y tế trở thành một chủ đề đáng quan tâm. Không chỉ dừng lại ở việc khám bệnh, khả năng điều khiển robot từ xa được mở rộng cho các ứng dụng tiềm năng như robot trợ giúp phẫu thuật từ xa, kiểm tra sức khỏe từ xa trước và sau khi điều trị, hoặc dùng để huấn luyện kĩ thuật phẫu thuật từ xa \cite{Su2022MixedRealityEnhancedHI}. Trong không gian vũ trụ, dự án phi hành gia robot của NASA được phát triển là một robot điều khiển từ xa \cite{zang2023FiveDimensional} dùng cho công việc bảo trì trạm không gian thay cho con người. Một ứng dụng khác của robot thay con người làm việc ở môi trường rủi ro cao là các robot điều khiển từ xa được vận hành để xử lý và dọn dẹp khu vực sau vụ nổ hạt nhân ở Fukushima, Nhật Bản vào năm 2011, \cite{Kawatsuma2012EmergencyResponseBy}. Việc đưa vào sử dụng các robot đã giúp giảm thiểu đáng kể chi phí và rủi ro điều động con người vào khu vực nhiễm phóng xạ nghiêm trọng. 

Có thể nói, mô hình cộng tác Người - Máy vẫn là một lĩnh vực đóng vai trò hữu ích để giải quyết các công việc có tính chất nguy hiểm cao đối với con người, cũng như là một giải pháp để giảm thiểu thời gian chậm trễ gây ra khi không yêu cầu các chuyên gia có phải mặt trực tiếp mà có thể thông qua phương thức teleoperation - hệ thống Robot và thực hiện công việc chuyên môn mang tính phức tạp cao. Mặc dù các công nghệ Trí tuệ nhân tạo (AI) và Học máy đang phát triển, các hệ thống robot hiện tại vẫn gặp khó khăn trong việc tự chủ giải quyết các tình huống bất ngờ nằm ngoài dữ liệu huấn luyện. Do đó, phương pháp Điều khiển từ xa (Teleoperation) tiếp tục là một giải pháp thiết thực, cho phép chuyên gia con người giám sát và trực tiếp vận hành hệ thống robot từ một khoảng cách an toàn \cite{darvish2023teleoperation}.

\section{Ứng dụng Thực tế ảo (VR) trong Điều khiển từ xa}
Sự thành công của một hệ thống điều khiển từ xa phụ thuộc phần lớn vào Giao diện tương tác Người - Máy (Human-Robot Interface). Giao diện cần phải đáp ứng hai nhu cầu: 1. Giao diện phải dễ nhìn, dễ sử dụng và 2. độ trễ điều khiển phải thấp để đảm bảo cảm giác điều khiển ổn định. Quan sát các hệ thống điều khiển công nghiệp truyền thống, người vận hành thường phải sử dụng màn hình 2D, bàn phím, tay cầm (Joystick) hoặc bộ điều khiển cầm tay (Teach Pendant) để tương tác với robot \cite{hetrick2020comparing}. Tuy nhiên, các cơ chế điều khiển này thường không trực quan, dẫn đến thời gian học sử dụng lâu dài và yêu cầu người có kinh nghiệm vận hành lâu dài mới có thể sử dụng hiệu quả. Sử dụng phần mềm qua màn hình 2D, bàn phím hoặc chuột yêu cầu người dùng liên tục quan sát không gian qua camera 2D, đối chiếu với mô hình 3D trong tư duy và thực hiện các thao tác cơ học. Điều này vô hình trung tạo ra tải nhận thức (cognitive load) lớn, dễ gây mệt mỏi và làm giảm độ chính xác khi thực hiện các tác vụ phức tạp \cite{chen2023comparing}. Ngoài ra, việc thiếu sự phản hồi không gian ba chiều khiến cho phương thức điều khiển truyền thống khó hiểu đối với người mới, yêu cầu thời gian huấn luyện và tập sử dụng trước khi được phép vận hành.

Để khắc phục những hạn chế của giao diện 2D nói riêng và phương thức điều khiển truyền thống nói chung, công nghệ Thực tế ảo (Virtual Reality - VR) đã được nghiên cứu và ứng dụng như một giải pháp thay thế hiệu quả. Công nghệ VR được xem là đóng vai trò quan trọng trong việc phát triển hệ thống cộng tác Người - Máy hiện đại. VR có thể được sử dụng để giải quyết hai thách thức của việc xây dựng ứng dụng điều khiển từ xa hiệu quả, đó là cải thiện khả năng điều khiển dễ dàng, cùng với phản hồi thị giác trực quan, cung cấp nhận thức về môi trường xung quanh thiết bị được điều khiển cho người sử dụng \cite{Luu2025EnhancingRealTime}. Không gian VR cung cấp khả năng hiển thị chiều sâu tự nhiên, cho phép người dùng quan sát môi trường làm việc của robot dưới dạng 3D thông qua kính hiển thị đội đầu (HMD) \cite{arclab2025beavr}. Trong một số nghiên cứu, việc điều khiển robotic arm - một ứng dụng thường thấy trong các nghiên cứu điều khiển từ xa bằng VR - có thể linh hoạt tuân theo chuyển động của bàn tay người, nhờ cơ chế ánh xạ phức tạp giữa tọa độ và góc xoay của đầu cánh tay robot với tọa độ và góc xoay của tay cầm dưới sự điều khiển của người vận hành \cite{Su2022MixedRealityEnhancedHI}. Thiết bị VR đã mở đường cho nhiều cách thức điều khiển đa dạng khác nhau, với mục tiêu đem lại cảm giác điều khiển trực quan, thân thiện và dễ hiểu nhất cho người dùng, qua đó cải thiện thời gian làm quen với ứng dụng điều khiển từ xa, giúp người mới có thể vận hành hiệu quả thiết bị robot trong thời gian ngắn, hoặc lý tưởng hơn là không cần trải qua huấn luyện. 

Lợi ích của việc sử dụng VR không chỉ dừng ở khả năng điều khiển dễ dàng, mà còn là việc cung cấp cảm quan về môi trường, vật thể và đối tượng điều khiển. Khi kết hợp với công nghệ Bản sao số (Digital Twin), người vận hành có thể có góc nhìn trực quan với một mô hình ảo của robot được đồng bộ trạng thái với robot thực tế, biết được phạm vi vận hành chính xác và cảm nhận chiều sâu mà không cần đến trực tiếp thiết bị. Phương pháp này hỗ trợ cải thiện nhận thức không gian (spatial awareness) và giúp các thao tác vận hành trở nên tự nhiên hơn \cite{torrejon2026teleoperation}. Đây chính là một trong những yếu tố khiến việc điều khiển qua ứng dụng VR có tiềm năng tốt hơn so với các ứng dụng điều khiển truyền thống, khi người dùng không còn phải tự ánh xạ giữa hình ảnh 2D/ 3D trên màn hình máy tính với mô hình 3D tư duy, giảm thiểu đáng kể tải nhận thức, giúp đem lại khả năng vận hành trong thời gian dài mà vẫn đảm bảo chất lượng công việc. 

\section{Sự nâng cấp về nền tảng điều khiển Robot thông dụng}

Để xây dựng một ứng dụng VR trực quan và đem lại khả năng điều khiển hiệu quả, không thể thiếu vai trò của một bộ điều khiển Robot hiệu năng cao. Ứng dụng VR sẽ không đóng vai trò trực tiếp điều khiển Robot, mà luôn thông qua một nền tảng vận hành robot để đảm bảo mọi hoạt động, hành vi của robot được điều phối ổn định và nhanh chóng. Về nền tảng vận hành robot, ROS1 (Robot Operating System) đã từ lâu đóng vai trò là một nền tảng tiêu chuẩn trong nghiên cứu robot, và là lựa chọn thường thấy trong hầu hết các nghiên cứu mở liên quan tới việc điều khiển nhiều loại robot khác nhau, từ robot xe (mobile robot), cánh tay robot (robotic arm), đến cả những robot phức tạp mô phỏng con người (humanoid robot). Tuy nhiên, kiến trúc của ROS1 bộc lộ một số hạn chế khi triển khai các hệ thống yêu cầu tính thời gian thực (real-time). ROS1 ban đầu được thiết kế chủ yếu dành cho mục đích nghiên cứu trong phòng thí nghiệm, nên nó thiếu đi những ràng buộc chặt chẽ, đặc biệt là về thời gian thực để có thể áp dụng trong môi trường sản xuất thực tế \cite{StevenRos2Usecases}. ROS1 sử dụng kiến trúc tập trung với một Master Node duy nhất và giao thức truyền thông TCPROS \cite{macenski2022robot}. Với việc tập trung mọi giao tiếp qua Roscore, việc giao tiếp mang rủi ro của "Single point of failure", khi node trung tâm gặp vấn đề, mọi giao tiếp trong hệ thống sẽ bị sụp đổ. Thứ hai, giao tiếp của ROS1 là TCP-ROS, dựa trên TCP/IP truyền thống, do đó chịu ảnh hưởng của vấn đề "Chặn đầu hàng đợi" (Head-of-line blocking). Đây là vấn đề khó tránh khỏi của giao thức dựa trên TCP/IP, do giao thức này yêu cầu dữ liệu truyền theo đúng thứ tự, nên khi trong luồng truyền dữ liệu có xảy ra tình trạng có những dữ liệu quan trọng, kích thước nhỏ và những dữ liệu lớn (như camera, point cloud) được truyền cùng lúc, nghẽn băng thông hoặc lỗi mạng ở những khối dữ liệu lớn sẽ "chặn" toàn bộ đường truyền, và các gói tin quan trọng cũng sẽ không truyền đi được, không thỏa mãn yêu cầu của hệ thống thời gian thực với độ ưu tiên. Vấn đề này xảy ra trong hệ thống điều khiển qua VR, với dữ liệu truyền tải thường bao gồm cả luồng hình ảnh camera (băng thông lớn) để cải thiện cảm nhận môi trường, cùng tín hiệu điều khiển/ trạng thái động học (kích thước nhỏ nhưng cần tần số cao) lại có độ ưu tiên cao để đảm bảo sự đồng bộ của Digital Twin. Việc truyền tải đồng thời qua kiến trúc của ROS1 dễ dẫn đến tình trạng nghẽn mạng cục bộ hoặc trễ gói tin (packet loss), từ đó làm giảm độ mượt mà của thao tác điều khiển \cite{kawshan2026digital, casini2025survey}. Thêm nữa, ROS1 không có cơ chế quản lý chất lượng dịch vụ truyền (QoS), do đó không có khái niệm về độ ưu tiên, khó đáp ứng với môi trường cần sự phản hồi nhanh hoặc khẩn cấp.

Nhận biết những giới hạn của ROS1, ROS2 đã được đề xuất và phát triển từ năm 2014, với mục tiêu cải tiến để trở thành một nền tảng vận hành robot quy mô công nghiệp. ROS2 thay thế giao thức TCP-ROS của mình và sử dụng Data Distribution Service (DDS) - một tiêu chuẩn công nghiệp mở về truyền thông dữ liệu đã được chứng minh trong các hệ thống không gian vũ trụ, quân sự và tài chính \cite{StevenRos2Usecases}. DDS hoạt động theo cơ chế mạng ngang hàng (peer-to-peer), cho phép các thiết bị tự động khám phá và giao tiếp trực tiếp với nhau mà không cần Master node, từ đó loại bỏ triệt để điểm lỗi duy nhất của ROS1. Đa số các cấu hình của DDS sử dụng giao thức UDP, không yêu cầu truyền lại các gói tin bị mất trừ khi được cấu hình, giúp tránh được độ trễ do Head-of-line blocking. Cơ chế xác định chất lượng dịch vụ (Quality of Service) cũng được xây dựng cho ROS2, qua đó cho phép người dùng định nghĩa các mức QoS khác nhau cho các luồng dữ liệu theo độ ưu tiên: Luồng dữ liệu lớn không cần real-time (như hình ảnh, video) có thể cấu hình là "Best Effort" - dữ liệu có thể được bỏ qua khi mạng trong tình trạng nghẽn, còn những dữ liệu nhỏ, quan trọng được cấu hình "Reliable" để luôn được ưu tiên trên đường truyền, đảm bảo việc điều khiển và cập nhật trạng thái luôn chính xác và đúng lúc. Nhờ sử dụng giao thức Data Distribution Service (DDS), ROS2 cung cấp cơ chế mạng ngang hàng phân tán và các cấu hình Chất lượng dịch vụ (QoS) linh hoạt. Điều này cho phép hệ thống phân luồng và ưu tiên các gói tin điều khiển quan trọng, góp phần duy trì tính ổn định của hệ thống trong môi trường truyền thông đa luồng \cite{macenski2022robot, ali2025digital}.

\section{Động lực và mục tiêu của đề tài}

Nhận biết nhu cầu phát triển một hệ thống cộng tác Người - Máy vẫn có vai trò thực tiễn và khả năng ứng dụng tốt, đề tài chủ trương xây dựng một hệ thống HRC thật sự hiệu quả, thông qua phương hướng nâng cấp từ kiểu điều khiển thông dụng qua phần mềm 2D, chuột hay bàn phím, lên phương thức điều khiển bằng ứng dụng VR với thiết bị đeo đầu là Meta Quest 3 - thiết bị VR/AR mới nhất hiện nay được phát hành bởi Meta, để dem lại một trải nghiệm điều khiển chân thực, dễ sử dụng và trực quan cho người dùng. Đối tượng điều khiển sẽ là một cánh tay robot công nghiệp Denso VS-6577 với sáu trục tự do, giao tiếp thông qua bộ điều khiển trung gian RC5. Không dừng ở đó, mục tiêu của đồ án hướng đến xây dựng một nền tảng điều khiển robot vững chắc với ROS2, tận dụng lợi thế về hạ tầng giao tiếp và cơ chế ưu tiên (QoS) của ROS2 để đảm bảo hiệu quả giao tiếp thời gian thực, độ trễ thấp, từ đó tăng cường chất lượng điều khiển và hiệu quả hoàn thành công việc trên ứng dụng VR giao tiếp với ROS2 so với các hình thức điều khiển truyền thống. Mục tiêu của đề tài nỗ lực xây dựng một nền tảng (framework) từ thấp tới cao, tích hợp những công nghệ, giải pháp hiện đại và tiêu chuẩn, tạo tiền đề để có thể dễ dàng phát triển và mở rộng sau này. Chi tiết các mục tiêu là:

\begin{enumerate}
    \item \textbf{Xây dựng cơ chế nhận lệnh và điều khiển trên RC5 để vận hành cánh tay robot VS-6577}: VS-6577 là một cánh tay robot công nghiệp, nên nó sẽ thường đi kèm với một bộ controller riêng biệt của hãng - RC5. Tuy nhiên, bộ điều khiển RC5 được phát triển với tiêu chí chương trình được tải và chạy theo chu trình và logic đã xác định - được viết sử dụng ngôn ngữ và phần mềm độc quyền của nhà sản xuất, và không được hỗ trợ để có thể chịu sự can thiệp của người vận hành từ xa. Chính vì vậy, chương trình vận hành Robot (ROS2) sẽ tập trung ở lớp cao hơn, trên RC5 sẽ chỉ tập trung nhận lệnh điều khiển và điều khiển trục robot vận hành ổn định trong khoảng thời gian dài, dưới sự giới hạn của cách thức giao tiếp duy nhất với module bên ngoài là giao tiếp UART.
    \item \textbf{Xây dựng chương trình vận hành chính bằng ROS2 với kiến trúc tiêu chuẩn}: ROS2 trở thành lựa chọn hiển nhiên cho đề tài, nhờ những lợi ích vượt trội so với ROS1 trong mục tiêu xây dựng ứng dụng thời gian thực để điều khiển từ xa cánh tay robot. Nhưng đề tài không dừng ở việc sử dụng ROS2 để giao tiếp - mà nhắm đến việc xây dựng một kiến trúc phần mềm điều khiển ROS2 tiêu chuẩn được khuyến khích, tích hợp đầy đủ ROS2-Control cùng Hardware Interface, MoveIt2, phân lớp rõ ràng các module và viết toàn bộ bằng ngôn ngữ C++ để đạt hiệu năng tốt nhất cho chương trình chạy thay vì Python như ROS1. Việc phân rõ các lớp thiết kế giúp cho hệ thống trở thành một platform với khả năng mở rộng thêm tính năng dễ dàng trong tương lai.
    \item \textbf{Xây dựng ứng dụng VR chạy độc lập trên kính Meta Quest 3}: Ứng dụng VR được xây dựng với Unity, chính là giao diện của hệ thống tương tác Người - Máy (HRC), cho phép người dùng có thể sử dụng để điều khiển cánh tay robot Denso VS-6577 từ xa. Tiêu chí xây dựng ứng dụng cần đảm bảo hai ý: 1. dễ sử dụng, dễ hiểu và đem lại nhiều phương thức điều khiển đa dạng, chính xác; 2. khả năng đồng bộ thời gian thực với cánh tay robot thực tế thông qua mô hình Digital Twin, tận dụng cơ chế ưu tiên dữ liệu quan trọng của ROS2 để luôn đảm bảo dữ liệu Digital Twin được cập nhật nhanh chóng, kịp thời, trong khi đó đem lại góc nhìn đa chiều, tổng quát về không gian hoạt động xung quanh thiết bị thực, đảm bảo độ hiệu quả sử dụng ứng dụng trong các công việc so với phương thức truyền thống.
\end{enumerate}

Dựa trên những mục tiêu này, quá trình phát triển của đồ án tới nay đã có một số kết quả, với những đóng góp nổi bật như:

\begin{itemize}
    \item \textbf{Đề xuất giải pháp điều khiển RC5 và cánh tay VS-6577 cho phép giao tiếp với module bên ngoài ổn định và tin cậy}: Mặc dù có giới hạn về phương thức giao tiếp với ROS2 (chỉ UART) - gồm giới hạn về tốc độ truyền và khả năng lỗi gói tin làm ngưng hệ thống, dự án đề xuất phương án phát hiện, xử lý và loại bỏ tác động ngừng hệ thống khi lỗi gói tin. Ngoài ra, cơ chế Ring Buffer cho phép tách biệt giữa luồng nhận dữ liệu và xử lý để xoay cánh tay robot, giúp cánh tay robot có thể xoay mượt khi có chuỗi lệnh điều khiển gửi xuống từ ROS2.
    \item \textbf{Hiện thực thành công kiến trúc điều khiển ROS2, tích hợp đầy đủ Ros2-Control và MoveIt}: Hiện thực ROS2 của dự án tuân thủ chính xác kiến trúc được đề xuất của Ros2-Control cho việc xây dựng một hệ thống điều khiển Robot có sự phân lớp rõ ràng, dễ dàng chỉnh sửa, nâng cấp và mở rộng, phân biệt giữa tầng Driver, tầng Hardware Interface và tầng Controllers, phía trên là lớp ứng dụng dựa trên MoveIt2, đem lại khả năng điều khiển và kiểm thử linh hoạt. Gazebo cũng được tích hợp vào hệ thống, đóng vai trò là một môi trường mô phỏng tích hợp chung với Ros2-Control, cho phép lớp ứng dụng có thể nhanh chóng phát triển và kiểm thử trước khi chạy trên robot thật, giảm thiểu rủi ro về mặt phần cứng. Đồng thời, đề tài cố gắng tối ưu về mặt thời gian giữa tần số điều khiển và hoạt động thực tế của cánh tay (dưới tính chất vật lý và ma sát) để có độ trễ giữa điều khiển và hoàn thành lệnh khoảng 2-3s.
    \item \textbf{Lựa chọn và tích hợp Ros-Sharp để giao tiếp ROS2 với thiết bị VR}: Ros-Sharp là một thư viện đang được tích cực phát triển bởi Siemens, để phổ biến hơn khả năng giao tiếp giữa ứng dụng chạy ROS1/ROS2 và ứng dụng Unity cho mô phỏng/ điều khiển. Ros-sharp hỗ trợ tốt ở cả phía ROS2 và Unity, cho phép tạo kết nối thông qua Web Socket, nhưng vẫn đảm bảo ứng dụng Unity trên thiết bị VR có thể tự do truy cập vào các topic, gọi Service, sử dụng Action tự do trong khi giữ được độ hiệu quả và thời gian thực của giao tiếp ROS2. Ros-Sharp được thiết kế khả mở rộng - khuyến khích lập trình viên thiết kế và chỉnh sửa các module, thêm các message/service/actions để đáp ứng nhu cầu của ứng dụng Unity-VR.
    \item \textbf{Phát triển ứng dụng VR với Unity trực quan và ổn định}: Ngoài Ros-Sharp, ứng dụng Unity cũng tích hợp gói phát triển Meta-SDK, để đem lại hiệu năng tốt và ổn định trên thiết bị Meta Quest 3, đồng thời thiết kế giao diện mang tính tương tác cao - cho phép người dùng được thoải mái hoạt động, thao tác trong môi trường ảo, trong khi đó cung cấp ba phương thức điều khiển cánh tay robot Denso VS-6577 kèm dữ liệu từ hai camera, đảm bảo người dùng có được cảm nhận về không gian của thiết bị được điều khiển. Thêm nữa, dữ liệu thực tế của cánh tay luôn được ưu tiên xử lý và cập nhật vào mô hình Digital Twin của cánh tay trong ứng dụng, cho phép người dùng dễ dàng nhận biết tình trạng, hình dáng và hoạt động của cánh tay trong không gian mà không gây tải nhận thức như các phương thức truyền thống. Ứng dụng hoạt động ổn định trong điều kiện bình thường, mức FPS đạt 72 trong nhiều lần thí nghiệm của dự án, đảm bảo trải nghiệm mượt mà và không chóng mặt, phù hợp để sử dụng trong thời gian dài. 
\end{itemize}

\section{Bố cục của báo cáo}

Phần tiếp theo của báo cáo sẽ bao gồm các phần như sau:

\begin{itemize}
    \item \textbf{Chương 2: Các nghiên cứu liên quan} Trình bày các nghiên cứu liên quan tới vấn đề xây dựng ứng dụng VR cho việc điều khiển Robot, thường đánh sâu vào giải quyết các vấn đề thường gặp trong hệ thống cộng tác Người - Máy trong việc điều khiển từ xa.
    \item \textbf{Chương 3: Cơ sở lý thuyết} Trình bày về những thành phần được sử dụng trong việc xây dựng hệ thống điều khiển của đồ án, giới thiệu về cánh tay Robot VS-6577, bộ điều khiển RC5, Bộ tay gắp được tích hợp (Gripper-A) gắn ở đầu cánh tay, sơ lược về ROS2, cùng Ros2-Control và MoveIt2 cho một kiến trúc điều khiển tiêu chuẩn, sau cùng là thiết vị kính VR Meta Quest 3, chạy ứng dụng VR được phát triển bởi game Engine Unity, tận dụng Ros-Sharp (Ros\#) cho giao tiếp hiệu quả và nhanh chóng giữa ứng dụng VR và ROS2.
    \item \textbf{Chương 4: Thiết kế và hiện thực hệ thống} Phần thiết kế hệ thống và cách thức đề tài hiện thực được trình bày chi tiết, từ tổng quát tới cụ thể ba phần chính: 1. chương trình điều khiển RC5 và VS-6577, 2. Thiết kế phần mềm ROS2 và 3. Thiết kế phần mềm điều khiển VR trên Unity Engine. 
    \item \textbf{Chương 5: Đánh giá kết quả} Trình bày kết quả hiện thực, và những thí nghiệm được sử dụng để đánh giá độ hiệu quả của module đã hiện thực. Đối với hệ thống điều khiển ROS2, trước nhất là độ hoàn thiện của hệ thống điều khiển, và khả năng vận hành thực tế phần cứng tới vị trí theo yêu cầu. Tiếp đến, độ trễ của một lần điều khiển MoveIt được đánh giá sau khi tối ưu về thời gian điều khiển và nhiều cấu hình liên quan tới robot trên cả ROS2 và RC5. Đối với ứng dụng VR, các tính năng kết nối, điều khiển đều hoạt động được, độ trễ giao tiếp giữa lệnh truyền từ ứng dụng VR tới ROS2 được kiểm tra thông qua xem xét RTT (Rount trip time) của gói tin đo đạc riêng. Sau cùng, dữ liệu hoạt động ứng dụng (CPU, RAM, FPS) được trích xuất sau phiên hoạt động dài, nhằm phân tích hiệu quả ứng dụng có phù hợp để đem lại trải nghiệm ổn định, không chóng mặt cho người sử dụng 
    \item \textbf{Chương 6: Tổng kết và hướng phát triển} Cuối cùng, chương 6 tổng kết lại những kết quả đã hoàn thành được của dự án, cũng như những điểm yếu cần cải thiện, đề xuất hướng phát triển của dự án sau này.
\end{itemize}